\chapter{Problem Formulation and Experimental Framework}
\label{ch:problem}

\section{Formal Problem Definition}

Let $D = \{x_1, x_2, \dots, x_N\}$ denote a collection of $N$ classical music recordings drawn from the iPalpiti archive.
Each item $x_i$ represents a complete audio track.

An embedding model is defined as a function:

\[
f_\theta : \mathcal{X} \rightarrow \mathbb{R}^d
\]

where $\theta$ denotes the model parameters and $d$ the embedding dimensionality.
The embedding of item $x_i$ is:

\[
z_i = f_\theta(x_i)
\]

Similarity between two items is computed using a similarity function $s(\cdot,\cdot)$, such as cosine similarity:

\[
s(z_q, z_i) = \frac{z_q \cdot z_i}{\|z_q\|\|z_i\|}
\]

Higher similarity values indicate greater embedding-space proximity.


Given a query item $x_q \in D$, all other items are ranked
in descending order of similarity $s(z_q, z_i)$.

Retrieval quality is evaluated with respect to a relevance function
$r_q(x_i) \in \{0,1\}$ defined by proxy tasks.

Ranking quality is measured using standard IR metrics including NDCG@K, Precision@K, and Recall@K.



\section{Definition of Embedding Capacity}

The primary independent variable in this study is embedding-model capacity.


Embedding capacity is operationalized primarily by model parameter count ($|\theta|$),
with embedding dimensionality ($d$) and architectural class treated as structural contributors.

For clarity, we consider capacity along three measurable aspects:

1. Parameter count ($|\theta|$)
2. Embedding dimensionality ($d$)
3. Architectural class (e.g., CNN-based vs. Transformer-based)

Models are categorized along a capacity spectrum ranging from lightweight convolutional architectures to larger transformer-based representations.

This operationalization allows capacity to be treated as a measurable experimental variable rather than a qualitative label (e.g., "heavy" vs. "light").


\section{Proxy Retrieval Tasks}

Relevance is defined via structured proxy tasks.

Let $g(x)$ denote a metadata-derived label function.

Two proxy families are defined:

\paragraph{Sanity Proxy.}
Relevance is defined by shared composer identity:
\[
r_q(x_i) = 1 \quad \text{if} \quad g_{\text{composer}}(x_q) = g_{\text{composer}}(x_i)
\]

\paragraph{Primary Musical Proxy.}
Relevance is defined by shared character labels (e.g., heavy vs. light):
\[
r_q(x_i) = 1 \quad \text{if} \quad g_{\text{character}}(x_q) = g_{\text{character}}(x_i)
\]

These proxy tasks are diagnostic instruments rather than end goals.
They allow controlled evaluation of embedding-space organization under constrained archival conditions.


\section{Experimental Variables and Controls}

This study follows a hypothesis-driven experimental design.

\paragraph{Independent Variable.}
Embedding-model capacity (Section 3.2).

\paragraph{Dependent Variables.}
Ranking-based retrieval metrics (NDCG@K, Precision@K, Recall@K).

\paragraph{Controlled Variables.}
\begin{itemize}
    \item Audio preprocessing pipeline
    \item Indexing mechanism
    \item Similarity function
    \item Evaluation protocol
\end{itemize}

By holding controlled variables constant, observed performance differences can be attributed to embedding capacity.

\section{Backend Pipeline Architecture}

The backend evaluation pipeline consists of four stages:

1. Audio preprocessing
2. Embedding extraction
3. Similarity indexing
4. Ranking and evaluation

All components other than embedding extraction are held constant across experiments.
The pipeline isolates embedding capacity as the central variable of interest.


\section{Conditioning Assumptions}

All claims in this thesis are conditioned on:

\begin{itemize}
    \item Small-scale dataset size ($N$ limited)
    \item Single-domain classical music archive
    \item Retrieval-based proxy evaluation
    \item Absence of user interaction modeling
\end{itemize}

Results should therefore be interpreted within these constraints.

These constraints define the validity domain of the study.
Findings should not be generalized beyond similar small-scale, single-domain archival settings without further empirical validation.




